\documentclass[12pt]{article}
\usepackage{amsmath, amssymb, geometry, enumitem, xcolor}
\geometry{letterpaper, margin=1in}

\title{Linear Algebra: Solution Sets of Linear Systems \\ \large Problem Sets and Solutions}
\author{}
\date{}

\begin{document}

\maketitle
\tableofcontents
\newpage

% ---------------------------------------------------------
\section{Concept Checklist \& Question Mapping}
% ---------------------------------------------------------
To ensure complete coverage of all concepts, tricks, and problem archetypes presented in the chapter, the following checklist maps each core concept to its corresponding problem(s) in the problem set.

\begin{itemize}
    \item \textbf{Existence of Nontrivial Solutions (Homogeneous):} Determining if $A\mathbf{x} = \mathbf{0}$ has free variables. \\ \textit{Questions: 1, 3, 20(a)}
    \item \textbf{The "Wide Matrix" Cheat:} Matrices with more columns than rows ($n > m$) automatically have nontrivial solutions. \\ \textit{Questions: 2}
    \item \textbf{Parametric Vector Form (Homogeneous):} Translating RREF of $A\mathbf{x} = \mathbf{0}$ into a vector span. \\ \textit{Questions: 4, 5, 6, 7}
    \item \textbf{Parametric Vector Form (Nonhomogeneous):} Solving $A\mathbf{x} = \mathbf{b}$ and expressing as $\mathbf{x} = \mathbf{p} + \mathbf{v}_h$. \\ \textit{Questions: 8, 9, 10}
    \item \textbf{Geometric Translation ($\mathbf{p} + \mathbf{v}_h$):} Understanding how a nonhomogeneous solution relates geometrically to its homogeneous counterpart. \\ \textit{Questions: 11, 29}
    \item \textbf{Geometric Constructions (Lines):} Deriving parametric vector forms using two given points ($\mathbf{x} = \mathbf{p} + t(\mathbf{q} - \mathbf{p})$). \\ \textit{Questions: 12, 13}
    \item \textbf{Reverse Engineering Matrices:} Constructing matrices that possess a specific null space. \\ \textit{Questions: 14, 15}
    \item \textbf{Creating Inconsistency via Column Space:} Choosing vectors $\mathbf{b}$ that break the pattern of the matrix columns. \\ \textit{Questions: 16, 17, 18}
    \item \textbf{Conceptual Synthesis (True/False):} Testing the definitions of consistency, pivots, and linearity. \\ \textit{Questions: 19, 20, 21, 22}
    \item \textbf{Variation 1: The Rank-Nullity Connection:} Using Rank + Nullity = $n$ to find the number of free variables without row reducing. \\ \textit{Questions: 23, 24}
    \item \textbf{Variation 2: Geometric Intersection of Hyperplanes:} Understanding dimension drops when adding independent equations to a system. \\ \textit{Questions: 25, 26}
    \item \textbf{Advanced Diagnostic ("Find the Flaw"):} Identifying subtle algebraic and conceptual errors. \\ \textit{Questions: 27, 28, 29, 30}
\end{itemize}

\newpage

% ---------------------------------------------------------
\section{Problem Sets}
% ---------------------------------------------------------

\subsection*{Set 1: Existence of Nontrivial Solutions}
\textbf{1.} Determine if the following homogeneous system has a nontrivial solution. Use as few row operations as possible.
\begin{align*}
    x_1 - 2x_2 + 3x_3 &= 0 \\
    2x_1 - 4x_2 + 6x_3 &= 0 \\
    -1x_1 + 2x_2 - 3x_3 &= 0
\end{align*}

\textbf{2.} Without performing any row operations, explain why the homogeneous system $A\mathbf{x} = \mathbf{0}$ must have a nontrivial solution if $A$ is given by:
$$A = \begin{bmatrix} 1 & -5 & 2 & 7 \\ -2 & 3 & 4 & 1 \end{bmatrix}$$

\textbf{3.} Let $A = \begin{bmatrix} 1 & 2 \\ 0 & 1 \\ 0 & 0 \end{bmatrix}$. Does $A\mathbf{x} = \mathbf{0}$ have a nontrivial solution? Explain geometrically what the solution set looks like.

\subsection*{Set 2: Parametric Vector Form (Homogeneous Systems)}
\textbf{4.} Write the solution set of the following homogeneous system in parametric vector form:
\begin{align*}
    x_1 + 2x_2 - x_3 &= 0 \\
    2x_1 + 3x_2 + 4x_3 &= 0
\end{align*}

\textbf{5.} Write the solution set of the given homogeneous system in parametric vector form:
\begin{align*}
    x_1 + x_2 - x_3 - x_4 &= 0 \\
    x_1 - x_2 - 3x_3 + 7x_4 &= 0 \\
    2x_1 - 4x_3 + 6x_4 &= 0
\end{align*}

\textbf{6.} Describe the solutions of $A\mathbf{x} = \mathbf{0}$ in parametric vector form, where $A$ is row equivalent to the given RREF matrix:
$$\begin{bmatrix} 1 & -3 & 0 & 2 & 0 \\ 0 & 0 & 1 & -4 & 0 \\ 0 & 0 & 0 & 0 & 1 \end{bmatrix}$$

\textbf{7.} A $3 \times 5$ homogeneous system $A\mathbf{x} = \mathbf{0}$ reduces to the following RREF matrix. Express its solution in parametric vector form.
$$\begin{bmatrix} 1 & 2 & 0 & -1 & 5 \\ 0 & 0 & 1 & 3 & -2 \\ 0 & 0 & 0 & 0 & 0 \end{bmatrix}$$

\subsection*{Set 3: Parametric Vector Form (Nonhomogeneous Systems)}
\textbf{8.} Write the solution set of the following nonhomogeneous system in parametric vector form:
\begin{align*}
    x_1 + 2x_2 - x_3 &= 5 \\
    2x_1 + 3x_2 + 4x_3 &= 9
\end{align*}

\textbf{9.} Describe the solutions of the following nonhomogeneous system in parametric vector form:
\begin{align*}
    x_1 + x_2 - x_3 - x_4 &= 4 \\
    x_1 - x_2 - 3x_3 + 7x_4 &= -2 \\
    2x_1 - 4x_3 + 6x_4 &= 2
\end{align*}

\textbf{10.} Attempt to find the parametric vector form of the solution set for the system below. What happens, and what does this mean geometrically?
\begin{align*}
    x_1 + x_2 + x_3 &= 1 \\
    2x_1 + 2x_2 + 2x_3 &= 3
\end{align*}

\textbf{11.} Give a geometric description of the solution set from Problem 8 and compare it to the solution set of the homogeneous system in Problem 4. 

\subsection*{Set 4: Geometric Constructions}
\textbf{12.} Find a parametric equation of the line $M$ passing through vectors $\mathbf{p} = \begin{bmatrix} 1 \\ -2 \\ 4 \end{bmatrix}$ and $\mathbf{q} = \begin{bmatrix} 3 \\ 0 \\ -1 \end{bmatrix}$.

\textbf{13.} Find a parametric equation of the line in $\mathbb{R}^4$ passing through $\mathbf{p} = \begin{bmatrix} 0 \\ 5 \\ -3 \\ 2 \end{bmatrix}$ and $\mathbf{q} = \begin{bmatrix} 1 \\ 1 \\ 1 \\ 1 \end{bmatrix}$.

\textbf{14.} Construct a $2 \times 2$ matrix $A$ such that the solution set of $A\mathbf{x} = \mathbf{0}$ is the line in $\mathbb{R}^2$ passing through $(2, 3)$ and the origin.

\textbf{15.} Construct a $3 \times 3$ matrix $B$ (with no zero rows) such that the solution set of $B\mathbf{x} = \mathbf{0}$ is the line spanned by the vector $\mathbf{v} = \begin{bmatrix} 1 \\ 1 \\ -1 \end{bmatrix}$.

\subsection*{Set 5: Creating Inconsistency}
\textbf{16.} Let $A = \begin{bmatrix} 3 & -2 \\ -6 & 4 \end{bmatrix}$. Find a vector $\mathbf{b}$ in $\mathbb{R}^2$ such that the solution set of $A\mathbf{x} = \mathbf{b}$ is empty.

\textbf{17.} Let $C = \begin{bmatrix} 1 & -2 & 3 \\ 1 & -2 & 3 \\ 1 & -2 & 3 \end{bmatrix}$. Describe the condition that a vector $\mathbf{b} = \begin{bmatrix} b_1 \\ b_2 \\ b_3 \end{bmatrix}$ must satisfy for $C\mathbf{x} = \mathbf{b}$ to be consistent.

\textbf{18.} True or False: If we find a vector $\mathbf{b}$ such that $A\mathbf{x} = \mathbf{b}$ is inconsistent, we can still use Theorem 6 ($\mathbf{x} = \mathbf{p} + \mathbf{v}_h$) to describe its solution set by setting $\mathbf{p} = \mathbf{0}$. Explain.

\subsection*{Set 6: Conceptual Synthesis and True/False}
\textbf{19.} True or False: A homogeneous system of equations can sometimes be inconsistent if the coefficient matrix has a row of zeros. Explain.

\textbf{20.} Let $A$ be a $4 \times 4$ matrix with four pivot positions. 
(a) Does the equation $A\mathbf{x} = \mathbf{0}$ have a nontrivial solution? 
(b) Does the equation $A\mathbf{x} = \mathbf{b}$ have a solution for every possible $\mathbf{b}$?

\textbf{21.} True or False: If a $3 \times 5$ matrix $A$ has 3 pivot rows, then $A\mathbf{x} = \mathbf{b}$ has a unique solution for every vector $\mathbf{b}$ in $\mathbb{R}^3$.

\textbf{22.} True or False: If $\mathbf{u}$ and $\mathbf{v}$ are both solutions to a nonhomogeneous system $A\mathbf{x} = \mathbf{b}$ (where $\mathbf{b} \neq \mathbf{0}$), then their sum $(\mathbf{u} + \mathbf{v})$ is also a solution to $A\mathbf{x} = \mathbf{b}$.

\subsection*{Set 7: Untested Variations (Rank-Nullity and Hyperplanes)}
\textbf{23.} Without solving, determine how many parameters (free variables) are needed to describe the solution set of a homogeneous system with a $4 \times 7$ coefficient matrix that has exactly 3 pivot columns.

\textbf{24.} An $8 \times 6$ matrix $A$ has Rank($A$) = 4. How many free variables does the system $A\mathbf{x} = \mathbf{0}$ have?

\textbf{25.} Consider a homogeneous linear system $A\mathbf{x} = \mathbf{0}$ containing 3 linearly independent equations with 5 variables. Geometrically, what is the dimension of the solution space?

\textbf{26.} Geometrically describe the solution set of a single linear equation $2x_1 + 3x_2 - x_3 + x_4 = 0$ in $\mathbb{R}^4$. 

\subsection*{Set 8: Advanced Diagnostic Testing ("Find the Flaw")}
\textit{For problems 27-30, act as the grader. Identify the subtle mathematical error, explain why it is wrong in one sentence, and provide the correct final answer.}

\textbf{27. Flawed Solution A:}
\textbf{Task:} Write the solution to $2x_1 - x_2 = 0$ in parametric vector form.
\textbf{Flawed Solution:} We isolate $x_1$ to get $x_1 = 0.5x_2$. Since $x_2$ is defined by $2x_1$, we set up our vector:
$$ \mathbf{x} = \begin{bmatrix} x_1 \\ x_2 \end{bmatrix} = \begin{bmatrix} x_1 \\ 2x_1 \end{bmatrix} = x_1 \begin{bmatrix} 1 \\ 2 \end{bmatrix} $$

\textbf{28. Flawed Solution B:}
\textbf{Task:} A system $A\mathbf{x} = \mathbf{b}$ reduces to the equations $x_1 = 4 - 2x_3$ and $x_2 = -1 + 3x_3$. Write the solution in parametric vector form.
\textbf{Flawed Solution:} We isolate the free variable $x_3$ and combine the terms:
$$ \mathbf{x} = \begin{bmatrix} x_1 \\ x_2 \\ x_3 \end{bmatrix} = x_3 \begin{bmatrix} 4 - 2 \\ -1 + 3 \\ 1 \end{bmatrix} = x_3 \begin{bmatrix} 2 \\ 2 \\ 1 \end{bmatrix} $$

\textbf{29. Flawed Solution C:}
\textbf{Task:} Describe geometrically the solution set of $\mathbf{x} = \begin{bmatrix} 3 \\ 0 \\ 2 \end{bmatrix} + t \begin{bmatrix} -1 \\ 1 \\ 0 \end{bmatrix}$.
\textbf{Flawed Solution:} Because the parameter $t$ can be any real number (including zero), this equation geometrically represents a line that passes through the origin in $\mathbb{R}^3$.

\textbf{30. Flawed Solution D:}
\textbf{Task:} Construct a $2 \times 2$ matrix $A$ such that the system $A\mathbf{x} = \mathbf{0}$ has NO trivial solution.
\textbf{Flawed Solution:} To prevent the trivial solution $\mathbf{x} = \mathbf{0}$, we must make the matrix inconsistent by having a row of zeros. Let $A = \begin{bmatrix} 1 & 1 \\ 0 & 0 \end{bmatrix}$. This forces inconsistency, ensuring no trivial solution exists.

\newpage

% ---------------------------------------------------------
\section{Detailed Solutions}
% ---------------------------------------------------------

\textbf{Problem 1:}
\textbf{Step 1: Form coefficient matrix.}
$A = \begin{bmatrix} 1 & -2 & 3 \\ 2 & -4 & 6 \\ -1 & 2 & -3 \end{bmatrix}$
\textbf{Step 2: Row operations.}
Replace $R_2$ with $R_2 - 2R_1$: $\begin{bmatrix} 1 & -2 & 3 \\ 0 & 0 & 0 \\ -1 & 2 & -3 \end{bmatrix}$
Replace $R_3$ with $R_3 + R_1$: $\begin{bmatrix} 1 & -2 & 3 \\ 0 & 0 & 0 \\ 0 & 0 & 0 \end{bmatrix}$
\textbf{Conclusion:} Column 1 has a pivot. Columns 2 and 3 lack pivots (free variables). Because free variables exist, the system has a nontrivial solution.
\textbf{Final Answer:} Yes, nontrivial solutions exist.

\textbf{Problem 2:}
\textbf{Explanation:} The matrix $A$ is $2 \times 4$. The maximum number of pivots any matrix can have is limited by its smallest dimension (here, 2 rows means at most 2 pivots). Since there are 4 columns (variables) and at most 2 pivots, there must be at least $4 - 2 = 2$ free variables. By the "Wide Matrix Cheat", a system with free variables always has nontrivial solutions.
\textbf{Final Answer:} Since $n > m$ (4 columns > 2 rows), free variables are mathematically guaranteed, yielding nontrivial solutions.

\textbf{Problem 3:}
\textbf{Explanation:} The matrix $A = \begin{bmatrix} 1 & 2 \\ 0 & 1 \\ 0 & 0 \end{bmatrix}$ is already in echelon form. Column 1 has a pivot (row 1) and Column 2 has a pivot (row 2). There are 0 free variables. Therefore, the only solution is the trivial solution $\mathbf{x} = \mathbf{0}$.
\textbf{Final Answer:} No nontrivial solution. Geometrically, the solution set is just a single point: the origin in $\mathbb{R}^2$.

\textbf{Problem 4:}
\textbf{Step 1: Row Reduce augmented matrix.}
$\begin{bmatrix} 1 & 2 & -1 & 0 \\ 2 & 3 & 4 & 0 \end{bmatrix} \xrightarrow{R_2 - 2R_1} \begin{bmatrix} 1 & 2 & -1 & 0 \\ 0 & -1 & 6 & 0 \end{bmatrix} \xrightarrow{-R_2} \begin{bmatrix} 1 & 2 & -1 & 0 \\ 0 & 1 & -6 & 0 \end{bmatrix}$
$\xrightarrow{R_1 - 2R_2} \begin{bmatrix} 1 & 0 & 11 & 0 \\ 0 & 1 & -6 & 0 \end{bmatrix}$
\textbf{Step 2: Extract equations.}
$x_1 + 11x_3 = 0 \implies x_1 = -11x_3$
$x_2 - 6x_3 = 0 \implies x_2 = 6x_3$
$x_3$ is free.
\textbf{Final Answer:} $\mathbf{x} = x_3 \begin{bmatrix} -11 \\ 6 \\ 1 \end{bmatrix}$

\textbf{Problem 5:}
\textbf{Step 1: Row Reduce to RREF.}
$\begin{bmatrix} 1 & 1 & -1 & -1 \\ 1 & -1 & -3 & 7 \\ 2 & 0 & -4 & 6 \end{bmatrix} \xrightarrow[R_3 - 2R_1]{R_2 - R_1} \begin{bmatrix} 1 & 1 & -1 & -1 \\ 0 & -2 & -2 & 8 \\ 0 & -2 & -2 & 8 \end{bmatrix}$
$\xrightarrow{R_3 - R_2} \begin{bmatrix} 1 & 1 & -1 & -1 \\ 0 & -2 & -2 & 8 \\ 0 & 0 & 0 & 0 \end{bmatrix} \xrightarrow{-\frac{1}{2}R_2} \begin{bmatrix} 1 & 1 & -1 & -1 \\ 0 & 1 & 1 & -4 \\ 0 & 0 & 0 & 0 \end{bmatrix}$
$\xrightarrow{R_1 - R_2} \begin{bmatrix} 1 & 0 & -2 & 3 \\ 0 & 1 & 1 & -4 \\ 0 & 0 & 0 & 0 \end{bmatrix}$
\textbf{Step 2: Extract equations.}
$x_1 = 2x_3 - 3x_4$, $x_2 = -x_3 + 4x_4$. $x_3, x_4$ are free.
\textbf{Final Answer:} $\mathbf{x} = x_3 \begin{bmatrix} 2 \\ -1 \\ 1 \\ 0 \end{bmatrix} + x_4 \begin{bmatrix} -3 \\ 4 \\ 0 \\ 1 \end{bmatrix}$

\textbf{Problem 6:}
\textbf{Step 1: Identify basic and free variables.}
Basic: $x_1, x_3, x_5$. Free: $x_2, x_4$.
\textbf{Step 2: Equations.}
$x_1 = 3x_2 - 2x_4$, $x_3 = 4x_4$, $x_5 = 0$.
\textbf{Final Answer:} $\mathbf{x} = x_2 \begin{bmatrix} 3 \\ 1 \\ 0 \\ 0 \\ 0 \end{bmatrix} + x_4 \begin{bmatrix} -2 \\ 0 \\ 4 \\ 1 \\ 0 \end{bmatrix}$

\textbf{Problem 7:}
\textbf{Step 1: Extract equations from RREF.}
$x_1 = -2x_2 + x_4 - 5x_5$, $x_3 = -3x_4 + 2x_5$. Free: $x_2, x_4, x_5$.
\textbf{Final Answer:} $\mathbf{x} = x_2 \begin{bmatrix} -2 \\ 1 \\ 0 \\ 0 \\ 0 \end{bmatrix} + x_4 \begin{bmatrix} 1 \\ 0 \\ -3 \\ 1 \\ 0 \end{bmatrix} + x_5 \begin{bmatrix} -5 \\ 0 \\ 2 \\ 0 \\ 1 \end{bmatrix}$

\textbf{Problem 8:}
\textbf{Step 1: Row reduce augmented matrix.}
$\begin{bmatrix} 1 & 2 & -1 & 5 \\ 2 & 3 & 4 & 9 \end{bmatrix} \xrightarrow{\text{RREF}} \begin{bmatrix} 1 & 0 & 11 & 3 \\ 0 & 1 & -6 & 1 \end{bmatrix}$
\textbf{Step 2: Equations.}
$x_1 = 3 - 11x_3$, $x_2 = 1 + 6x_3$. $x_3$ is free.
\textbf{Final Answer:} $\mathbf{x} = \begin{bmatrix} 3 \\ 1 \\ 0 \end{bmatrix} + x_3 \begin{bmatrix} -11 \\ 6 \\ 1 \end{bmatrix}$

\textbf{Problem 9:}
\textbf{Step 1: Row Reduce augmented matrix.} (Uses steps from Prob 5)
$\begin{bmatrix} 1 & 1 & -1 & -1 & 4 \\ 1 & -1 & -3 & 7 & -2 \\ 2 & 0 & -4 & 6 & 2 \end{bmatrix} \xrightarrow{\text{RREF}} \begin{bmatrix} 1 & 0 & -2 & 3 & 1 \\ 0 & 1 & 1 & -4 & 3 \\ 0 & 0 & 0 & 0 & 0 \end{bmatrix}$
\textbf{Step 2: Equations.}
$x_1 = 1 + 2x_3 - 3x_4$, $x_2 = 3 - x_3 + 4x_4$.
\textbf{Final Answer:} $\mathbf{x} = \begin{bmatrix} 1 \\ 3 \\ 0 \\ 0 \end{bmatrix} + x_3 \begin{bmatrix} 2 \\ -1 \\ 1 \\ 0 \end{bmatrix} + x_4 \begin{bmatrix} -3 \\ 4 \\ 0 \\ 1 \end{bmatrix}$

\textbf{Problem 10:}
\textbf{Step 1: Reduce.}
$\begin{bmatrix} 1 & 1 & 1 & 1 \\ 2 & 2 & 2 & 3 \end{bmatrix} \xrightarrow{R_2 - 2R_1} \begin{bmatrix} 1 & 1 & 1 & 1 \\ 0 & 0 & 0 & 1 \end{bmatrix}$
\textbf{Conclusion:} The last row equates $0 = 1$, which is impossible.
\textbf{Final Answer:} The system is inconsistent (empty solution set). Geometrically, this represents two parallel planes that never intersect.

\textbf{Problem 11:}
\textbf{Comparison:} The solution in Problem 4 is a line through the origin in $\mathbb{R}^3$. The solution in Problem 8 is the exact same line, but it has been shifted (translated) away from the origin by the particular solution vector $\mathbf{p} = [3, 1, 0]^T$.
\textbf{Final Answer:} It is a line parallel to the homogeneous solution set, translated by $\mathbf{p} = [3, 1, 0]^T$.

\textbf{Problem 12:}
\textbf{Step 1: Find direction vector.} $\mathbf{v} = \mathbf{q} - \mathbf{p} = \begin{bmatrix} 3 \\ 0 \\ -1 \end{bmatrix} - \begin{bmatrix} 1 \\ -2 \\ 4 \end{bmatrix} = \begin{bmatrix} 2 \\ 2 \\ -5 \end{bmatrix}$.
\textbf{Step 2: Apply formula.} $\mathbf{x} = \mathbf{p} + t\mathbf{v}$.
\textbf{Final Answer:} $\mathbf{x} = \begin{bmatrix} 1 \\ -2 \\ 4 \end{bmatrix} + t \begin{bmatrix} 2 \\ 2 \\ -5 \end{bmatrix}$

\textbf{Problem 13:}
\textbf{Step 1: Direction vector.} $\mathbf{v} = \mathbf{q} - \mathbf{p} = \begin{bmatrix} 1 \\ 1 \\ 1 \\ 1 \end{bmatrix} - \begin{bmatrix} 0 \\ 5 \\ -3 \\ 2 \end{bmatrix} = \begin{bmatrix} 1 \\ -4 \\ 4 \\ -1 \end{bmatrix}$.
\textbf{Final Answer:} $\mathbf{x} = \begin{bmatrix} 0 \\ 5 \\ -3 \\ 2 \end{bmatrix} + t \begin{bmatrix} 1 \\ -4 \\ 4 \\ -1 \end{bmatrix}$

\textbf{Problem 14:}
\textbf{Step 1: Setup.} We need $A\mathbf{x} = \mathbf{0}$ for $\mathbf{x} = \begin{bmatrix} 2 \\ 3 \end{bmatrix}$. Let $A = \begin{bmatrix} a & b \\ c & d \end{bmatrix}$.
$2a + 3b = 0 \implies b = -\frac{2}{3}a$.
\textbf{Step 2: Pick values.} Let $a = 3$, so $b = -2$. Repeat for bottom row.
\textbf{Final Answer:} $A = \begin{bmatrix} 3 & -2 \\ 3 & -2 \end{bmatrix}$ (Other multiples work).

\textbf{Problem 15:}
\textbf{Step 1: Setup.} We need $B\mathbf{v} = \mathbf{0}$, so $b_{i1}(1) + b_{i2}(1) + b_{i3}(-1) = 0 \implies b_{i1} + b_{i2} = b_{i3}$ for each row $i$.
\textbf{Step 2: Choose independent rows.}
Row 1: $1, 0, 1$. Row 2: $0, 1, 1$. Row 3: $1, 1, 2$.
\textbf{Final Answer:} $B = \begin{bmatrix} 1 & 0 & 1 \\ 0 & 1 & 1 \\ 1 & 1 & 2 \end{bmatrix}$

\textbf{Problem 16:}
\textbf{Step 1: Analyze columns.} Col 2 is $-2/3$ of Col 1. The column space is spanned by $\begin{bmatrix} 3 \\ -6 \end{bmatrix}$ (or $\begin{bmatrix} 1 \\ -2 \end{bmatrix}$).
\textbf{Step 2: Pick an external vector.} Choose any vector where the bottom component is NOT $-2$ times the top.
\textbf{Final Answer:} $\mathbf{b} = \begin{bmatrix} 1 \\ 1 \end{bmatrix}$ guarantees inconsistency.

\textbf{Problem 17:}
\textbf{Step 1: Analyze column space.} All columns of $C$ are identical. Thus, any output vector $C\mathbf{x}$ will just be a scalar multiple of $\begin{bmatrix} 1 \\ 1 \\ 1 \end{bmatrix}$.
\textbf{Final Answer:} To be consistent, $\mathbf{b}$ must have identical components: $b_1 = b_2 = b_3$.

\textbf{Problem 18:}
\textbf{Explanation:} Theorem 6 explicitly requires the nonhomogeneous system to be consistent (meaning at least one particular solution $\mathbf{p}$ exists). If it is inconsistent, there is no $\mathbf{p}$. Setting $\mathbf{p} = \mathbf{0}$ would just yield the homogeneous solution, which maps to $\mathbf{0}$, not $\mathbf{b}$.
\textbf{Final Answer:} False. Theorem 6's premise requires consistency.

\textbf{Problem 19:}
\textbf{Explanation:} By definition, $A\mathbf{0} = \mathbf{0}$ is always true. The trivial solution is always a guaranteed solution regardless of row reductions.
\textbf{Final Answer:} False. Homogeneous systems are always consistent.

\textbf{Problem 20:}
\textbf{Step 1: Part a.} $4 \times 4$ with 4 pivots means 0 free variables. Only the trivial solution exists.
\textbf{Step 2: Part b.} 4 pivot rows means no row of zeros in RREF. Thus, it can never equal $[0, 0, 0, 0 \mid nonzero]$. It is consistent for all $\mathbf{b}$.
\textbf{Final Answer:} (a) No. (b) Yes.

\textbf{Problem 21:}
\textbf{Step 1: Consistency check.} 3 pivots in 3 rows means consistent for all $\mathbf{b}$.
\textbf{Step 2: Uniqueness check.} 5 columns but only 3 pivots means $5 - 3 = 2$ free variables. Free variables mean infinitely many solutions.
\textbf{Final Answer:} False. It has a solution for every $\mathbf{b}$, but it is not unique.

\textbf{Problem 22:}
\textbf{Step 1: Check linearity.} $A(\mathbf{u} + \mathbf{v}) = A\mathbf{u} + A\mathbf{v} = \mathbf{b} + \mathbf{b} = 2\mathbf{b}$. Because $2\mathbf{b} \neq \mathbf{b}$ (since $\mathbf{b} \neq \mathbf{0}$), the sum is not a solution.
\textbf{Final Answer:} False.

\textbf{Problem 23:}
\textbf{Step 1: Rank-Nullity.} Total columns $n = 7$. Rank (pivots) = 3. 
Nullity = $n - \text{Rank} = 7 - 3 = 4$.
\textbf{Final Answer:} Exactly 4 parameters.

\textbf{Problem 24:}
\textbf{Step 1: Rank-Nullity.} Total columns $n = 6$. Rank = 4.
$6 - 4 = 2$.
\textbf{Final Answer:} 2 free variables.

\textbf{Problem 25:}
\textbf{Step 1: Dimension drops.} We start in $\mathbb{R}^5$ (5 variables). Each linearly independent equation is a hyperplane that drops the dimension by 1. $5 - 3 = 2$.
\textbf{Final Answer:} The solution space is 2-dimensional (a plane through the origin).

\textbf{Problem 26:}
\textbf{Step 1: Variables.} $n = 4$, rank = 1. Free variables = 3.
\textbf{Final Answer:} Geometrically, it is a 3-dimensional volume (a hyperplane) passing through the origin in $\mathbb{R}^4$.

\textbf{Problem 27:}
\textbf{Flaw Identified:} The student used the basic variable $x_1$ as the free parameter.
\textbf{Explanation:} In RREF, parametric vectors must be spanned by the non-pivot free variables (in this case, $x_2$), not the basic variables.
\textbf{Correct Answer:} $x_1 = 0.5x_2 \implies \mathbf{x} = x_2 \begin{bmatrix} 0.5 \\ 1 \end{bmatrix}$.

\textbf{Problem 28:}
\textbf{Flaw Identified:} The student incorrectly factored out $x_3$ from the constant terms in the vector.
\textbf{Explanation:} Constants and variable terms must be separated into a particular vector $\mathbf{p}$ and a homogeneous vector $\mathbf{v}_h$; you cannot combine constants with a scaled parameter.
\textbf{Correct Answer:} $\mathbf{x} = \begin{bmatrix} 4 \\ -1 \\ 0 \end{bmatrix} + x_3 \begin{bmatrix} -2 \\ 3 \\ 1 \end{bmatrix}$.

\textbf{Problem 29:}
\textbf{Flaw Identified:} The student claimed the line passes through the origin.
\textbf{Explanation:} Because of the translation vector $\mathbf{p} = [3, 0, 2]^T$, which is not a multiple of the direction vector, the line is shifted away from the origin.
\textbf{Correct Answer:} It geometrically represents a line in $\mathbb{R}^3$ translated away from the origin by the vector $[3, 0, 2]^T$.

\textbf{Problem 30:}
\textbf{Flaw Identified:} The student assumed an inconsistent matrix has no trivial solution.
\textbf{Explanation:} A homogeneous system $A\mathbf{x} = \mathbf{0}$ can never be inconsistent; the trivial solution $\mathbf{x} = \mathbf{0}$ is always mathematically guaranteed regardless of matrix structure.
\textbf{Correct Answer:} It is impossible to construct a matrix where $A\mathbf{x} = \mathbf{0}$ lacks the trivial solution.

\newpage

% ---------------------------------------------------------
\section{Concept Checklist \& Question Mapping (Repeated)}
% ---------------------------------------------------------
To ensure complete coverage of all concepts, tricks, and problem archetypes presented in the chapter, the following checklist maps each core concept to its corresponding problem(s) in the problem set.

\begin{itemize}
    \item \textbf{Existence of Nontrivial Solutions (Homogeneous):} Determining if $A\mathbf{x} = \mathbf{0}$ has free variables. \\ \textit{Questions: 1, 3, 20(a)}
    \item \textbf{The "Wide Matrix" Cheat:} Matrices with more columns than rows ($n > m$) automatically have nontrivial solutions. \\ \textit{Questions: 2}
    \item \textbf{Parametric Vector Form (Homogeneous):} Translating RREF of $A\mathbf{x} = \mathbf{0}$ into a vector span. \\ \textit{Questions: 4, 5, 6, 7}
    \item \textbf{Parametric Vector Form (Nonhomogeneous):} Solving $A\mathbf{x} = \mathbf{b}$ and expressing as $\mathbf{x} = \mathbf{p} + \mathbf{v}_h$. \\ \textit{Questions: 8, 9, 10}
    \item \textbf{Geometric Translation ($\mathbf{p} + \mathbf{v}_h$):} Understanding how a nonhomogeneous solution relates geometrically to its homogeneous counterpart. \\ \textit{Questions: 11, 29}
    \item \textbf{Geometric Constructions (Lines):} Deriving parametric vector forms using two given points ($\mathbf{x} = \mathbf{p} + t(\mathbf{q} - \mathbf{p})$). \\ \textit{Questions: 12, 13}
    \item \textbf{Reverse Engineering Matrices:} Constructing matrices that possess a specific null space. \\ \textit{Questions: 14, 15}
    \item \textbf{Creating Inconsistency via Column Space:} Choosing vectors $\mathbf{b}$ that break the pattern of the matrix columns. \\ \textit{Questions: 16, 17, 18}
    \item \textbf{Conceptual Synthesis (True/False):} Testing the definitions of consistency, pivots, and linearity. \\ \textit{Questions: 19, 20, 21, 22}
    \item \textbf{Variation 1: The Rank-Nullity Connection:} Using Rank + Nullity = $n$ to find the number of free variables without row reducing. \\ \textit{Questions: 23, 24}
    \item \textbf{Variation 2: Geometric Intersection of Hyperplanes:} Understanding dimension drops when adding independent equations to a system. \\ \textit{Questions: 25, 26}
    \item \textbf{Advanced Diagnostic ("Find the Flaw"):} Identifying subtle algebraic and conceptual errors. \\ \textit{Questions: 27, 28, 29, 30}
\end{itemize}

\end{document}